\documentclass[11pt]{article}
\usepackage[utf8]{inputenc}
\usepackage[T1]{fontenc}
\usepackage{amsmath}
\usepackage{multicol}
\usepackage{geometry}
\usepackage{tikz}
\usetikzlibrary{shapes.geometric, arrows.meta}
\usepackage{enumitem}
\usepackage{xcolor}
\usepackage{titlesec}

% Configurações de layout
\geometry{a4paper, left=1cm, right=1cm, top=0.5cm, bottom=1.2cm}
\setlength{\columnseprule}{0.4pt}
\setlength{\baselineskip}{1.0\baselineskip}

% Cores para títulos
\titleformat{\section}{\normalfont\Large\bfseries\color{blue}}{\thesection}{1em}{}
\titleformat{\subsection}{\normalfont\large\bfseries\color{red}}{\thesubsection}{1em}{}
\titleformat{\subsubsection}{\normalfont\normalsize\bfseries\color{black}}{\thesubsubsection}{1em}{}

\title{\textcolor{blue}{Revisão: Perímetro e Área}}
\author{Professor: Jefferson}
\date{}

\begin{document}

\maketitle
\vspace{-1cm}

\begin{center}
\large{Nome: \underline{\hspace{8cm}} \quad Turma: \underline{\hspace{3cm}}}
\end{center}

\begin{multicols}{2}

\section*{1. Conceitos Fundamentais}

\subsection*{Perímetro}
Soma das medidas de todos os lados de uma figura plana.

\subsection*{Área}
Medida da superfície ocupada pela figura.

\section*{2. Fórmulas Essenciais}

\begin{itemize}
\item \textbf{Retângulo:} $P = 2(b+h)$; $A = b \cdot h$
\item \textbf{Quadrado:} $P = 4L$; $A = L^2$
\item \textbf{Triângulo:} $P = a+b+c$; $A = \frac{b \cdot h}{2}$
\item \textbf{Círculo:} $C = 2\pi r$; $A = \pi r^2$
\end{itemize}

\section*{3. Questões Resolvidas}

\begin{enumerate}

\item O perímetro de um retângulo é 50 cm, e seu comprimento é o dobro da largura. Qual é a largura?  
{\color{red} $2(x+2x)=50 \Rightarrow 6x=50 \Rightarrow x=8,33$ cm (largura).}

\item Um quadrado tem área igual a 144 m². Qual é o seu perímetro?  
{\color{red} $L=\sqrt{144}=12 \Rightarrow P=4\cdot 12=48$ m.}

\item O perímetro de um triângulo é 72 cm. Dois lados medem 20 cm e 18 cm. Qual é a medida do terceiro lado?  
{\color{red} $72-(20+18)=34$ cm.}

\item Um terreno retangular tem 5 m a mais de comprimento do que de largura. Se seu perímetro é 50 m, qual é a largura?  
{\color{red} $2(x+x+5)=50 \Rightarrow 4x+10=50 \Rightarrow x=10$ m.}

\item Um campo quadrado tem perímetro 60 m. Qual é a área?  
{\color{red} $L=60/4=15 \Rightarrow A=15^2=225$ m².}

\item A base de um retângulo mede 3 vezes sua altura. Se a área é 75 m², qual é a altura?  
{\color{red} $3x \cdot x=75 \Rightarrow 3x^2=75 \Rightarrow x=5$ m (altura).}

\item Um losango tem perímetro de 64 cm. Qual é a medida de cada lado?  
{\color{red} $64/4=16$ cm.}

\item Um retângulo tem 2 m a mais de largura que de comprimento. Se a área é 120 m², qual é o comprimento?  
{\color{red} $x(x+2)=120 \Rightarrow x^2+2x-120=0 \Rightarrow x=10$ m (comprimento).}

\item O perímetro de um triângulo equilátero é 45 cm. Qual é a área?  
{\color{red} $L=45/3=15 \Rightarrow A=\frac{\sqrt{3}}{4}(15^2)\approx 97,4$ cm².}

\item Um retângulo tem comprimento igual ao triplo da largura. Se seu perímetro é 64 cm, qual é a área?  
{\color{red} $2(x+3x)=64 \Rightarrow 8x=64 \Rightarrow x=8$ cm, $A=8 \cdot 24=192$ cm².}

\item Um círculo tem área 78,5 m² (use $\pi = 3,14$). Qual seu perímetro?  
{\color{red} $r^2=78,5/3,14=25 \Rightarrow r=5 \Rightarrow C=2\pi r=31,4$ m.}

\item Um trapézio tem bases 8 cm e 12 cm, e altura 5 cm. Qual sua área?  
{\color{red} $A=\frac{(8+12)\cdot 5}{2}=50$ cm².}

\item Um hexágono regular tem lado 5 cm. Qual seu perímetro?  
{\color{red} $6\cdot 5=30$ cm.}

\item Um triângulo retângulo tem catetos 6 cm e 8 cm. Qual seu perímetro?  
{\color{red} Hipotenusa $=10$, então $P=6+8+10=24$ cm.}

\item Um paralelogramo tem lados 7 cm e 9 cm. Qual seu perímetro?  
{\color{red} $2(7+9)=32$ cm.}

\item Um retângulo tem área 48 cm². Se aumentarmos seus lados em 25\%, qual a nova área?  
{\color{red} $A_{novo}=48 \cdot (1,25)^2=75$ cm².}

\item Um jardim quadrado tem área 81 m². Se reduzirmos seu lado em 2 m, qual o novo perímetro?  
{\color{red} $L=9 \Rightarrow L'=7 \Rightarrow P'=4\cdot 7=28$ m.}

\item Um retângulo tem perímetro 30 cm e área 50 cm². Quais suas dimensões?  
{\color{red} $x+y=15$, $xy=50 \Rightarrow$ soluções: $x=10$, $y=5$.}

\item Um triângulo isósceles tem perímetro 36 cm e base 14 cm. Qual é a medida dos lados iguais?  
{\color{red} $2x+14=36 \Rightarrow x=11$ cm.}

\item Um terreno retangular tem razão entre lados 3:4 e perímetro 84 m. Qual sua área?  
{\color{red} $2(3k+4k)=84 \Rightarrow 14k=84 \Rightarrow k=6$, lados $=18$ e $24$, área $=432$ m².}

\end{enumerate}

\end{multicols}

\end{document}

